\chapter{Testing the model behind the figure}\label{ch:falsification}

Many statisticians hold a reasonable belief about risk models fitted inside an
institution. The model that produces a disclosure risk figure is fitted to the
institution's own register: its key frequencies, its error rates, its records.
Checking that model against the same register looks circular. The fit was chosen
to agree with those data, so a test on those data can only confirm the fit, and
any real check must wait for an outside attack that nobody hopes to see. On this
view the risk model is an assumption to be argued about, and the register has
nothing more to say.

The belief fails because a linkage model does more than describe the register. It
predicts how often an attack will name the right record, and it predicts this
separately for each level of confidence the adversary reaches. Those predictions
can be compared with attack outcomes that the fit never used. On the simulated
population of Chapter~\ref{ch:simulation}, two linkage models were run side by side.
One used the adversary's true error precision. The other assumed the adversary's
covariate was noisier than it was. The two risk figures differed by 0.007. % source: paper/generated/numbers.tex (runattackcovariate 0.4642, runmisspecifiedrate 0.4571)
Their stated posteriors, set against the realised success rates on held-out
attacks, differed by a factor of seven: a mean gap of 0.007 for the correct model
and 0.049 for the wrong one. % source: paper/generated/numbers.tex (runcalibrationgapcorrect, runcalibrationgapwrong)
The register did contradict one of the models. The risk figure could not see it,
and a calibration test could.

This chapter asks one question. What evidence can show that the model behind a
risk figure is wrong, and what should the institution do with the figure when it
is?

\section{An implementation check is not a validation}

The pipeline that produces this book's numbers reproduces every equation of the
earlier chapters in code, and a test suite confirms that the code computes what the
equations say. That is an implementation check. It shows that the reciprocal of a
key frequency is computed correctly when Assumptions~\ref{as:key} to~\ref{as:unif}
of Chapter~\ref{ch:scenario} hold. It cannot show that they hold. A model can be
coded without error and still describe an adversary nobody faces, a population
that does not exist, or a mechanism that behaves differently from the way it was
described.

Empirical validation needs a different kind of evidence. A test of the model must
name a quantity that the model predicts, a source of evidence the model was not
fitted on, a band inside which agreement counts as agreement, and a consequence if
the evidence falls outside it. All four are written down before the evidence is
read. A test without a stated band is a description, and a test without a
consequence is a report that nobody has to act on.

At the Bank, Naledi keeps a folder for each assessment, and the first sheet in it
is the list of tests she will run, with their bands, signed and dated before she
starts the simulations for Tomás's request. The chair once asked why the sheet
had to come first. Her answer was that a band chosen after the result is read can
always be drawn to include the result.

The long form of the paper behind this book proposes six tests. Figure~\ref{c15:fig-six} sets them out by the part
of the model each one examines and by what follows from a failure. Four of them,
linkage calibration, population-frequency transport, adversary stress and
mechanism monotonicity, test the model of a single release, and a failure of any
of them makes the release-level figure invalid. The fifth, the composition test,
tests whether the release is the right object of assessment at all. The sixth,
temporal transport, cannot be run when the assessment is made, and the expiry date
of the governance threshold enforces it instead.

\begin{figure}[ht]
\centering
\begin{tikzpicture}[font=\small,
  tst/.style={draw,align=center,text width=3.1cm,minimum height=1.05cm},
  inv/.style={tst,fill=gray!20,pattern=north east lines,pattern color=gray!60},
  prg/.style={tst,pattern=dots,pattern color=gray!70},
  exp/.style={tst,dashed}]
  \node[inv] (cal) at (0,0) {1. Linkage calibration\\\scriptsize the linkage model};
  \node[inv] (frq) at (4.2,0) {2. Frequency transport\\\scriptsize the population model};
  \node[inv] (str) at (8.4,0) {3. Adversary stress\\\scriptsize the declared adversary};
  \node[inv] (mon) at (0,-1.7) {4. Monotonicity\\\scriptsize the mechanism};
  \node[prg] (cmp) at (4.2,-1.7) {5. Composition\\\scriptsize the object assessed};
  \node[exp] (tmp) at (8.4,-1.7) {6. Temporal transport\\\scriptsize the time of use};
  \node[draw,fill=gray!20,pattern=north east lines,pattern color=gray!60,minimum width=0.5cm,minimum height=0.35cm] at (-1.3,-3.2) {};
  \node[anchor=west,font=\footnotesize] at (-0.95,-3.2) {failure: figure invalid, no decision};
  \node[draw,pattern=dots,pattern color=gray!70,minimum width=0.5cm,minimum height=0.35cm] at (4.5,-3.2) {};
  \node[anchor=west,font=\footnotesize] at (4.85,-3.2) {failure: govern the programme};
  \node[draw,dashed,minimum width=0.5cm,minimum height=0.35cm] at (-1.3,-3.75) {};
  \node[anchor=west,font=\footnotesize] at (-0.95,-3.75) {cannot run at assessment: enforced by the expiry date};
\end{tikzpicture}
\caption{The six falsification tests, each with the part of the model it examines
(small type) and the consequence of a failure. Grey fill with hatching marks the
four tests whose failure makes a release-level figure invalid; dotted fill marks
the composition test, whose failure moves the assessment to the programme; the
dashed outline marks the test that waits for the passage of time.}
\label{c15:fig-six}
\end{figure}

\section{Linkage calibration}

The first test examines the linkage model itself. When the declared adversary
attacks a target, the model gives a posterior probability that the record it names
is the right one. If the model is correct, then among all attacks in which the
adversary named a record with posterior near 0.3, about 30 per cent should be
right. The test bins the attacks by stated posterior, computes the realised success
rate in each bin from held-out outcomes, and reports the mean gap between the two.

Figure~\ref{c15:fig-calib} shows the test on the simulated population. With the
adversary's error model stated correctly, the eight bins lie close to the diagonal
and the mean gap is 0.007. When the model assumes that the adversary's covariate is
2.5 times noisier than it is, the stated posteriors fall below the realised rates
in the middle of the range: in the bin from 0.3 to 0.5, the model says 0.399 and the
attacks succeed 0.496 of the time. % source: artifacts-paper/results.json attack.calibration (bin 0.3-0.5, misspecified); configs/paper.json misspecified_covariate_factor 2.5
The mean gap is 0.049. % source: paper/generated/numbers.tex runcalibrationgapwrong
Meanwhile the success rate of the attack hardly moves: 0.464 with the right model
and 0.457 with the wrong one. % source: paper/generated/numbers.tex runattackcovariate, runmisspecifiedrate

\begin{figure}[ht]
\centering
\includegraphics[width=0.58\textwidth]{fig3_calibration.pdf}
\caption{Linkage calibration on the simulated population: mean stated posterior
against realised success rate in eight bins, 20{,}000 attacks per curve. Circles on
a solid line show the declared model; triangles on a dashed line show a model that
overstates the adversary's covariate noise by a factor of 2.5. The dotted diagonal
is perfect calibration.}
\label{c15:fig-calib}
\end{figure}

The near equality of the two success rates is not an accident of this simulation.
Proposition~\ref{c15:prop-invariance} shows why a wrongly scaled error model leaves
the adversary's choice, and so its success rate, unchanged, while it changes the
confidence the adversary attaches to that choice.

\begin{proposition}[A misstated error scale leaves success unchanged]\label{c15:prop-invariance}
Let the target $t$ lie in a candidate set $C$ of released records, all equally
likely a priori in the adversary's information, with published values $x_j$ for
$j \in C$. The adversary observes $z = x_t + e$ with $e \sim N(0,\sigma^2)$, assumes a
standard deviation $s > 0$, forms the posterior $p_s(j) \propto
\exp\{-(z-x_j)^2/(2s^2)\}$ and names $\hat\jmath$, the candidate of largest posterior,
with ties broken by a fixed rule. Then
\begin{enumerate}
\item[(i)] $\hat\jmath$ does not depend on $s$, so the success probability
$\Prob(\hat\jmath = t)$ is the same for every $s > 0$;
\item[(ii)] when $s = \sigma$ the stated posterior is calibrated:
$\Prob(\hat\jmath = t \mid z) = p_\sigma(\hat\jmath)$;
\item[(iii)] the stated posterior $p_s(\hat\jmath)$ is non-increasing in $s$,
tends to $1/|C|$ as $s \to \infty$, and, when the nearest candidate is
unique, tends to 1 as $s \to 0$.
\end{enumerate}
\end{proposition}

Part (i) says that the risk figure, which counts how often the named record is
right, cannot reveal a wrong error scale. Parts (ii) and (iii) say where the
evidence lies instead: in the posteriors, which are right only under the right
scale. An adversary who believes its covariate is noisier than it is reports
posteriors too close to $1/|C|$, as the dashed curve in Figure~\ref{c15:fig-calib}
does through most of its range. The attack in the pipeline also carries error in
the key fields, so the candidate sets differ slightly between the two scales, and
the success rates differ by 0.007 instead of by nothing. The conclusion is the same.

This answers the belief stated at the start of the chapter. A model calibrated on
the register can be refuted by the register, because the model makes a family of
predictions and the fit pins down only some of them. The calibration test reads
the predictions the risk figure ignores.

The test has a weakness in the repository as it stands, and the reader should know
it. The configuration that drives the paper's run sets the band for the mean
calibration gap at 0.06. % source: configs/paper.json calibration_ece_limit 0.06
The misspecified model's gap of 0.049 lies inside that band, so the pipeline as
configured would have passed the wrong model. The pipeline also applies the test to
the declared model only, so the comparison in Figure~\ref{c15:fig-calib} is reported
and not gated. % source: src/privacy_risk/pipeline.py validate() uses attack["calibration"][0]
The long form describes the misspecified model as failing the test, which is true
of the sevenfold contrast and not of the configured band. A band has to be set from
the size of error that would change a decision, and a band of 0.06 on a gap whose
correct value is 0.007 is too wide to catch an error of the kind this chapter
studies. The band is the institution's choice, and this one should be tightened.

Naledi's sheet for Tomás's request sets the band at 0.02 and records why: on the
simulated population a correctly specified model sits well below it, and the
misstated scale sits well above it. The chair asked whether 0.02 was a standard
figure. It is not. It is Naledi's reading of what separates the two cases, signed
before she runs the attacks, and the reviewer can disagree with it in writing.

\section{Population-frequency transport}

The second test examines the model of the population. Journalist risk, the chance
that a sample-unique record is also unique in the population
(Chapter~\ref{ch:measures}), cannot be computed from the sample alone. It rests on a
fitted law for the population cell sizes, here a negative binomial, and the fitted
law makes predictions about quantities that can be checked against held-out or
external counts.

On the simulated population, where the truth is known, the fitted law recovers the
mean population cell size closely: 14.83 against a true 14.88 at a 5 per cent
sample. % source: paper/generated/numbers.tex runfittedmeancell 14.83, runtruemeancell 14.88
It recovers the uniqueness probability less well, with errors between 8 and 25 per
cent across five sampling fractions, and every error lies on the conservative side,
since the model's uniqueness probability exceeds the truth at each fraction.
% source: paper/generated/numbers.tex runjournalisterrorlow 8, runjournalisterrorhigh 25; artifacts-paper/results.json measures.journalist (model > truth at all five fractions)
The test also exposes a misfit that a single number hides. A negative binomial law
that is closed under random thinning would keep the same dispersion whatever the
sampling fraction. The fitted dispersion moves instead, from 0.33 at a 2 per cent
sample to 0.51 at 20 per cent and 0.59 at 50 per cent.
% source: artifacts-paper/results.json measures.journalist dispersion 0.325, 0.380, 0.441, 0.505, 0.594
The family does not describe this population, though it summarises it well enough
at one fraction.

The paper draws the consequence that its own test demands. A model-based
uniqueness probability that fails transport is replaced by a bound, or reported
beside the truth and not used. No decision in the long form rests on the
model-based figure, and on the public Census extract, where there is no truth to
check against, the uniqueness probability is reported across a range of sampling
fractions instead of as one value (Chapter~\ref{ch:public}).

Here too the code is narrower than the text. The pipeline's frequency test checks
the absolute error of the fitted mean cell size against a band of 1.5 cells.
% source: src/privacy_risk/validation.py frequency_transport_error; configs/paper.json frequency_mae_limit 1.5
The measured error is 0.055 cells, so the test passes, but a test of the mean cell
size would pass a model whose uniqueness probability was wrong by a factor of two.
% source: artifacts-paper/results.json validation.metrics.frequency_mae 0.0551
An institution that uses the journalist measure for a decision should test the
quantity it uses, the uniqueness probability or the share of small cells, against
external counts, and should set the band on that quantity.

\section{Adversary stress}

The third test examines the declared adversary. Every risk figure in this book is
conditional on what the adversary is declared to know, and the declaration is a
guess about sources outside the institution's control. The stress test asks
whether a materially stronger adversary that is still plausible would change the
decision. If it would, the release fails or the scenario is revised in writing; the
figure for the weaker adversary is not quietly kept.

The best-known demonstration that a declared adversary can be too weak is
\citet{sweeney2000}. Hospital discharge records released without names still
carried date of birth, sex and postcode, and a public voter list carried the same
fields with names attached. A release reviewed against an adversary who held only
the published file faced one who held a second file. The stress test is the
discipline of asking, before release, what the second file could be.

On the simulated population the stronger adversary holds a covariate twice as
precise as the declared one. % source: configs/paper.json stronger_covariate_factor 0.5
Across the 45 mechanism settings of the grid, it raises the simulated risk by between
1 and 76 per cent, with a median of 40 per cent.
% source: paper/generated/numbers.tex runstresslow 1, runstresshigh 76, runstressmedian 40; artifacts-paper/results.json validation.metrics.stress
The stronger adversary is never weaker at any setting, which is what the pipeline
checks: a capability that lowered the risk would point to a misdeclared scenario or
a coding error. The size of the movement is the more useful result. A setting that
passes the tolerance against the declared adversary with a margin of less than 40
per cent may well fail against an adversary with a better reading, and that is why
the long form takes the risk under the stronger adversary, the robust risk, as the
coordinate of the risk--utility frontier (Chapter~\ref{ch:tolerance}).

Capability and modelling are different things, and the test keeps them apart. An
adversary that keeps its information but misjudges its own error, the case of
Proposition~\ref{c15:prop-invariance}, succeeds as often as before, so the stress
test cannot see it; the calibration test can. An adversary with better information
succeeds more often, so the calibration test may pass while the stress test fails.
Neither test replaces the other.

The test also removed a protection lever from the pipeline. An early version
offered record swapping between districts. It lowered the simulated risk, but only
against an adversary that read the released district at face value, and nothing in
the declared scenario justified that assumption. An adversary who knew swapping
was used, or who held a second source for the district, gained the protection
back. The lever was taken out, and the paper reports the removal. % source: docs/CLAIMS_EVIDENCE.md "What was tried and dropped"
Reporting its protection would have turned an assumption about the adversary into
a property of the release.

When Naledi ran the stress test for Tomás's request, the stronger adversary was not
hypothetical. Pieter, a credit analyst at a reporting lender, knows the exposure of
his own bank's loans to the last unit. The declared outside adversary reads
exposure with error; Pieter does not. On the synthetic register an insider who reads
exposure with a log error of 0.01 and knows the loan's history succeeds against
77 per cent of targets with no suppression, and against 42 per cent at a minimum
cell count of 300. % source: paper/bis_ifc_full.tex tab:record insider row (0.772 at k=1, 0.416 at k=300)
No floor up to 300 stops him. Naledi's record therefore names him as a separate
adversary with his own row, governed by the access route and not by the floor.

\section{Mechanism monotonicity}

The fourth test examines the protection mechanism. When an institution claims that
a stronger setting reduces risk, a larger minimum cell count or a coarser key, the
claim is empirical and can fail. The test looks for settings where stronger
protection is followed by higher measured risk, and requires that each violation
be investigated and not smoothed away.

The long form reports four kinds of finding under this test. The first is a large,
systematic violation of a different claim: every row of the $k$-anonymity table
breaks the nominal bound $1/k$ on the worst target, with the worst target at
$k = 10$ identified with probability 0.744 against a nominal 0.1.
% source: paper/generated/table_kanonymity.tex (k=10 worst 0.744)
The investigation of that violation produced a proposition, that $k$-anonymity
delivers $1/k$ only when the adversary holds nothing beyond the key
(Chapter~\ref{ch:measures}). The other three are small breaks in order. The worst
target rises from 0.744 at $k = 10$ to 0.768 at $k = 20$, estimates that rest on 250
replicates each. % source: paper/generated/table_kanonymity.tex
The pooled rate on the credit register rises from 0.260 with no suppression to
0.269 under $k$-anonymity with $k = 5$, about one standard error at 4{,}000 attacks
each. % source: paper/generated/table_evaluation.tex
The large-firm identification rate rises from 0.247 to 0.273 between the same two
levels, on a few hundred targets. % source: paper/generated/table_credit.tex (k=1 0.247, k=5 0.273)
None of the three exceeds what sampling noise explains, and the long form treats
them as noise after looking for a mechanism and finding none.

The decision depends on this test less than it seems to, for a reason worth stating.
The choice of a floor in this book does not assume that risk falls as the floor
rises. Fixed-sequence testing (Chapter~\ref{ch:selection}) tests the floors from the
strictest down and stops at the first failure, and its error guarantee holds
whatever the order of the true risks. Monotonicity is tested because a claim that a
lever protects should be checked, not because the selection rule needs it.

The pipeline's version of the test runs over the minimum cell counts at the finest
key, with a band of 0.025 for simulation noise, and finds no violation on the paper's
run. % source: src/privacy_risk/pipeline.py validate(); configs/paper.json monotonicity_tolerance 0.025; artifacts-paper/results.json validation.metrics.monotonicity_violations []

\section{The composition test}

The fifth test asks whether the single release is the right thing to assess. It runs
the joint attack on the programme, compares the pooled risk with the interval the
marginal risks imply (Proposition~\ref{prop:boole}), and reports the pooled risk
against a programme tolerance of its own.

On the simulated population the test fails clearly. Two releases, each 3-anonymous,
pooled by an adversary who intersects their candidate sets, identify the target with
probability 0.411, which is 2.8 times the upper bound of 0.148 that their two
marginal risks imply. % source: paper/generated/numbers.tex runpooledtwo 0.4114; paper/bis_ifc_full.tex sec:falsify (2.8, bound 0.148)
Four releases reach 0.458, 1.8 times the bound of all four.
% source: paper/generated/numbers.tex runpooledall 0.4581, runpooledallratio 1.8

A failure of this test does not invalidate the release-level model. The linkage
model, the population model and the mechanism may all be right, and each release may
be assessed correctly on its own. What fails is the assumption that the release is
the object to govern. Proposition~\ref{prop:nopool} says that the marginal risks
cannot bound the pooled risk, so the consequence is a change of object: the
programme needs its own simulation, its own tolerance and its own horizon. The
pipeline records the composition result beside the validation report and keeps it
out of the validity flag for that reason. % source: src/privacy_risk/validation.py full_validation_report
It applies the result in the gate instead: on the paper's run, the chosen setting
passes on its own with an upper bound of 0.081, and the gate still returns a fail,
because the upper bound on the pooled risk of the declared four-release programme,
0.464, exceeds the release tolerance of 0.12. % source: artifacts-paper/decision.json (upper_risk 0.0806, pooled_upper_bound 0.4640, tolerance 0.12, status FAIL)

The decision record in the repository shows one loose end. The validation step
judges the pooled bound against a programme tolerance of 0.5, which it meets, while
the gate judges the same bound against the release tolerance of 0.12, which it does
not. % source: artifacts-paper/decision.json programme.programme_tolerance_used_by_validation 0.5; configs/paper.json programme.tau_programme 0.5
Both statements are in the file, and a reader has to work out which governed the
outcome. A record should carry one programme tolerance, owned and dated like the
release tolerance, and the gate should compare the pooled bound with it.

Tomás's request is a composition question before it is anything else. He asked for
a research extract in March, a second in June and a third with the year's vintages.
Each was assessed when it came in. Naledi's first act on the new request was to
assess the three together, with his auxiliary information held fixed across them,
because the earlier records did not license any statement about the set.

\section{Temporal transport}

The sixth test cannot be run when the assessment is made. An assessment is
calibrated on the register, the auxiliary sources and the linkage methods of its
date. New public sources appear, a register changes its coverage, and attack methods
improve. The test asks that the assessment fail or be re-estimated when any of these
moves observed attack performance outside the range the assessment declared.

Because the evidence for this test arrives later, the governance threshold enforces
it through the expiry date $d$ and the re-assessment triggers. The pipeline refuses
a decision on or after the expiry date. % source: src/privacy_risk/governance.py decide (assessed >= expiry -> FAIL)
A bound from Chapter~\ref{ch:monitoring} shows how little shift a certificate absorbs.
If the deployed attack distribution lies within a Kullback--Leibler divergence $\rho$
of the one assessed, the deployed risk is bounded through the Bernoulli divergence.
The programme floor of 150 on the credit register has a confirmed upper bound of
0.109, and it can absorb a divergence of about $6 \times 10^{-4}$ before that bound
reaches the tolerance of 0.12. % source: paper/bis_ifc_full.tex (kl(0.12||0.109) about 6e-4)
Halving the standard deviation of the adversary's exposure reading, from 0.6 to
0.3, is a divergence of about 0.32 per reading, three orders of magnitude larger,
and it moves the floor from 150 to a median of 300 over twenty register draws.
% source: paper/bis_ifc_full.tex (divergence 0.32 per reading); paper/generated/table_replication.tex
The certificate holds for the declared reading precision and for little else. The
expiry date exists because the world the certificate describes will move.

\section{Why a failed test blocks the decision}

When a test fails, the tempting response is to keep the figure and add a margin. If
the calibration gap is too large, raise the risk figure by a fifth and compare that
with the tolerance. The response is tempting because it is cheap and seems prudent.
It is wrong, and the reason is that a failed test does not say the figure is too
small by some amount. It says the model that produced the figure does not describe
the attack, and nothing in the evidence says by how much it errs.

Proposition~\ref{c15:prop-gate} makes both halves of the argument precise. A rule
that refuses to decide when a test fails can only remove passes, so it keeps the
error guarantee of the upper bound whatever the dependence between the test and the
bound. A fixed margin cannot repair a model that omits what the adversary knows,
because such a model can understate the risk by any factor.

\begin{proposition}[Blocking keeps the level; no margin repairs a wrong model]\label{c15:prop-gate}
Let $\overline{R}$ be an upper bound for the model risk $R$ with
$\Prob(\overline{R} < R) \le \alpha$ when the model is correct, let $V$ be any event,
possibly dependent on the same simulations, on which the validation tests pass, and
let the decision pass when $V$ occurs and $\overline{R} \le \tau$.
\begin{enumerate}
\item[(i)] If the model is correct and $R > \tau$, the decision passes with
probability at most $\alpha$.
\item[(ii)] For every margin $c \ge 1$, tolerance $\tau \in (0,1)$ and
$\varepsilon > 0$, there is a release and a model, satisfying the assumptions of
Proposition~\ref{prop:recip} for an adversary who holds only the key, such that
$c\,\overline{R} \le \tau$ with probability at least $1 - \varepsilon$, while an
adversary who also holds one published attribute exactly identifies every target
with probability one.
\end{enumerate}
\end{proposition}

Part (ii) is the $k$-anonymity finding in its plainest form. The model is correct
for the adversary it describes and wrong for the one who attacks, and no scaling of
its output closes the gap, because the gap depends on what the adversary holds and
not on the model's own error. The only repair is a new model, which means a new
declaration, new simulations and a new set of tests.

The pipeline encodes the rule. When a release-level test fails, the gate returns
\texttt{INVALID} and the process exits with a non-zero status, so a failed
validation cannot be presented as a release decision.
% source: src/privacy_risk/pipeline.py main(); src/privacy_risk/governance.py decide
The status differs from a fail. A fail says that the evidence, read by a model the
institution accepts, shows the risk above the tolerance; the remedies of
Chapter~\ref{ch:protocol} apply. An invalid result says the institution has no
accepted model and so no evidence either way. The release cannot go out on it, and
the work returns to the declaration.

The chair raised the margin idea at the committee meeting on Tomás's request. A
draft calibration check on a new auxiliary source had come back with a gap above
Naledi's band, and the chair proposed adding a fifth to the figure and moving on.
Naledi showed the committee the dashed curve of Figure~\ref{c15:fig-calib} and the
two success rates beside it. A larger figure would have looked prudent and would
have been computed from the same wrong model. The committee recorded the result as
invalid, and Naledi rebuilt the declaration.

\section{What the tests can and cannot show}

The six tests do not prove that a model is right. A model that passes all six can
still be wrong about an adversary nobody thought to declare, and no test can range
over adversaries that are not written down; Chapter~\ref{ch:open} returns to this
limit. What the tests do is give the model a stated chance to fail, so that a
figure that survives them has survived something. The claims-and-evidence register
of the repository states for each claim what would refute it, and the last row is
the governance protocol itself, whose claim to improve auditability rests on an
institutional argument and has no comparative evidence yet.
% source: docs/CLAIMS_EVIDENCE.md last row of the table

The answer to the chapter's question is short. Evidence can show that the model
behind a risk figure is wrong whenever the model predicts something beyond the
figure: posteriors that can be compared with outcomes, population counts that can
be compared with external benchmarks, an ordering of risk across settings, a
response to a stronger adversary, and a range of conditions over time. The
institution's own data can supply most of that evidence, so the belief that a model
fitted to the register cannot be refuted by it is false. When a test fails, the
figure is not adjusted. The assessment is invalid, no decision is taken on it, and
the model is rebuilt from a new declaration.

\begin{inpractice}[In practice: testing the model before the decision]
\setlength{\emergencystretch}{3em}%
Before running any attack, write the test sheet: the six tests, the quantity each
compares, the evidence it uses and the band for each, signed and dated. Set each
band from the size of error that would change the decision, not from habit. Test
the quantity the decision uses: if the journalist measure enters the decision, test
the uniqueness probability against external counts, not only the mean cell size.
Run a stronger plausible adversary for every scenario and name any insider as a
separate adversary. Treat a failed release-level test as an invalid assessment, not
as a high figure. Treat a failed composition test as a reason to assess the
programme. Put the expiry date and the re-assessment triggers in the record, since
they carry the sixth test.

\smallskip\noindent\textit{In the software.} The tests are in
\texttt{privacy\_\allowbreak{}risk.\allowbreak{}validation}: \texttt{calibration\_\allowbreak{}table} and
\texttt{calibration\_\allowbreak{}gap} for linkage calibration,
\texttt{frequency\_\allowbreak{}transport\_\allowbreak{}error} for the population model,
\texttt{stress\_\allowbreak{}ratios}, \texttt{monotonicity\_\allowbreak{}violations} and
\texttt{composition\_\allowbreak{}test}, combined by \texttt{full\_\allowbreak{}validation\_\allowbreak{}report}.
\texttt{privacy\_\allowbreak{}risk.\allowbreak{}pipeline.\allowbreak{}validate} runs them on a pipeline run, and
\texttt{governance.\allowbreak{}decide} returns \texttt{INVALID} when the report has not passed.
The bands come from the configuration file:
\begin{lstlisting}
report = full_validation_report(
    ece, float(config.get("calibration_ece_limit", 0.06)),
    transport, float(config.get("frequency_mae_limit", 1.5)),
    violations, stress, composition)
\end{lstlisting}
\end{inpractice}

\section{Proofs}

\begin{proof}[Proof of Proposition~\ref{c15:prop-invariance}]
(i) For every $s > 0$ the map $u \mapsto \exp\{-u^2/(2s^2)\}$ is strictly decreasing
in $u \ge 0$, and the prior is uniform over $C$. The posterior $p_s(j)$ is therefore
largest exactly for the candidates that minimise $|z - x_j|$. This set does not
depend on $s$, and the fixed tie-breaking rule picks the same member of it for
every $s$, so $\hat\jmath$ is a function of $z$ and the $x_j$ alone. The event
$\{\hat\jmath = t\}$ and its probability are then the same for every $s$.

(ii) When $s = \sigma$ the adversary's likelihood is the true conditional density of
$z$ given that $j$ is the target, and with a uniform prior on $C$ that matches the
target's distribution in the adversary's information, Bayes' theorem gives
$\Prob(t = j \mid z) = p_\sigma(j)$ for each $j \in C$. Since $\hat\jmath$ is a function
of $z$, $\Prob(\hat\jmath = t \mid z) = p_\sigma(\hat\jmath)$.

(iii) Write $d_j = (z - x_j)^2$ and $d_* = d_{\hat\jmath} = \min_j d_j$. Then
\[
p_s(\hat\jmath) = \Big(\sum_{j \in C} \exp\{-(d_j - d_*)/(2s^2)\}\Big)^{-1}.
\]
Each term with $d_j > d_*$ increases with $s$ and tends to 1 as $s \to \infty$, and
each term with $d_j = d_*$ equals 1, so $p_s(\hat\jmath)$ is non-increasing in $s$
and tends to $1/|C|$. If the nearest candidate is unique, every other term tends to
0 as $s \to 0$, so $p_s(\hat\jmath) \to 1$.
\end{proof}

\begin{proof}[Proof of Proposition~\ref{c15:prop-gate}]
(i) The decision passes on $V \cap \{\overline{R} \le \tau\}$, which is contained in
$\{\overline{R} \le \tau\}$. When $R > \tau$, $\{\overline{R} \le \tau\} \subseteq
\{\overline{R} < R\}$, whose probability is at most $\alpha$ when the model is
correct. No assumption on the dependence between $V$ and $\overline{R}$ is used.

(ii) Choose an integer $k \ge 2$ with $1/k < \tau/c$, and let $\delta = \tau/c - 1/k > 0$.
Build a release with $n = kM$ records, $M$ key classes of exactly $k$ records each,
and one further published attribute whose values are distinct across all $n$
records. For an adversary who holds only the key, Assumptions~\ref{as:key}
to~\ref{as:unif} hold and Proposition~\ref{prop:recip} gives risk $R = 1/k$ for every
target. Let $\overline{R}$ be the exact one-sided upper bound from $B$ simulated
attacks of that adversary. Its success count is binomial with mean $B/k$, and the
exact bound converges to $1/k$ in probability as $B \to \infty$, so for $B$ large
enough $\Prob(\overline{R} \le 1/k + \delta) \ge 1 - \varepsilon$, and on that event
$c\,\overline{R} \le c/k + c\delta = \tau$. An adversary who also holds the target's
value of the distinct attribute exactly finds a single released record with that
value, which is the target's, so it identifies every target with probability one.
\end{proof}
